\chapter{Đóng gói và cắt hai chiều}
\chapterlead{Trong đóng gói, mệnh đề không làm việc trực tiếp với hình chữ
nhật; nó làm việc với tọa độ, hướng đặt và các quan hệ tách. Chất lượng phép mã hóa
phụ thuộc vào việc chọn đúng giao diện giữa ba lớp đó.}

\index{strip packing}\index{bin packing}\index{non-overlap}
\index{rotation}\index{coordinate compression}

\flowdiagram{Vật và hướng đặt}{Tọa độ thứ tự}{Quan hệ tách}
  {Cận khung chứa}{Bố trí và chứng nhận}

\section{Bản đồ các phương pháp chính xác cho đóng gói}

Hai chiều tạo khó khăn khác một chiều: không chồng lấn là một phép tuyển giữa
hai trục, còn tọa độ có nhiều giá trị đối xứng trong vùng trống. Văn liệu về
phương pháp chính xác vì thế không xoay quanh một mô hình duy nhất. Khảo cứu
của Iori và cộng sự phân biệt các mô hình theo vị trí đặt, vị trí tương đối,
sinh tập/cột, lập trình ràng buộc và các kỹ thuật nhánh--cận chuyên biệt
\parencite{iori2021packing}.

Đối với SAT, ba họ biểu diễn là đối chứng trực tiếp:
\begin{description}[style=nextline]
  \item[Chiếm chỗ]
  Biến cho từng ô mà vật chiếm. Không chồng lấn trở thành cục bộ nhưng số biến
  phụ thuộc diện tích lưới; tính liên tục của một hình cần các mệnh đề liên kết.

  \item[Vị trí đặt tuyệt đối]
  Mỗi biến chọn một vị trí khả dĩ của vật. Đồ thị xung đột nối hai vị trí chồng
  lấn; bài toán trở thành chọn đúng một vị trí cho mỗi vật. Các phương pháp vị
  trí tuyệt đối và đóng gói hoàn hảo khai thác nén tọa độ, tổng tập con và suy
  luận trên không gian trống
  \parencite{huang2013rectangle}.

  \item[Vị trí tương đối + tọa độ]
  Mỗi vật có tọa độ; mỗi cặp có biến chứng trái/phải/trên/dưới. Mã hóa thứ tự
  nén cận tọa độ, còn mệnh đề quan hệ xử lý phép tuyển.
\end{description}

Ngoài SAT, lập trình ràng buộc dùng \textsc{NoOverlap}/diffn, các phép chiếu
tích lũy và suy luận năng lượng; MIP dùng hằng số lớn \(M\), bất đẳng thức
tuyển hoặc mẫu phủ tập. Các phương pháp chính xác cho đóng gói trực giao còn
kết hợp cận, quan hệ trội và nhánh--cận chuyên dụng
\parencite{clautiaux2007orthogonal,clautiaux2008cp}.

\begin{center}
\captionof{table}{Các biểu diễn chính xác cho đóng gói hai chiều.}
\label{tab:packing-formulations}
\begin{tabularx}{\textwidth}{@{}>{\bfseries\sffamily}p{.2\textwidth}YYY@{}}
\toprule
Biểu diễn & Quy mô theo & Lan truyền tự nhiên & Phù hợp khi\\
\midrule
Chiếm chỗ & Diện tích lưới & Xung đột ô & Khung chứa nhỏ/rất rời rạc\\
Vị trí đặt & Số vị trí khả dĩ & Đồ thị xung đột & Miền đã nén mạnh\\
Tọa độ--thứ tự & Tổng độ rộng miền & Cận và quan hệ tách & Miền nguyên vừa\\
CP toàn cục & Số vật + bộ lan truyền & Không chồng lấn/tích lũy & Cận hình học mạnh\\
MIP/phủ tập & Biến/cột vị trí & Thư giãn và nhát cắt & Cận LP hữu ích\\
\bottomrule
\end{tabularx}
\end{center}

Ba trục tạo nên đóng góp SAT cho đóng gói là nén miền tọa độ, làm mạnh lan
truyền không chồng lấn và tái sử dụng mệnh đề đã học giữa các cận. Kiến trúc
mới được nhận diện qua thay đổi có thể đo trên ít nhất một trong ba trục này.

\section{Mô hình hình học rời rạc}

Mỗi vật \(i\) là hình chữ nhật kích thước \(w_i\times h_i\). Tọa độ
\((x_i,y_i)\) chỉ góc trái dưới. Khi cho phép quay \(90^\circ\), biến
\(\rho_i\) xác định kích thước hiệu dụng
\[
  (\widehat w_i,\widehat h_i)=
  \begin{cases}
    (w_i,h_i), & \rho_i=0,\\
    (h_i,w_i), & \rho_i=1.
  \end{cases}
\]
Trong khung chứa \(W\times H\), điều kiện nằm gọn là
\[
  0\le x_i\le W-\widehat w_i,
  \qquad
  0\le y_i\le H-\widehat h_i.
\]

Hai vật \(i\) và \(j\) không giao nhau khi ít nhất một trong bốn quan hệ đúng:
\[
  x_i+\widehat w_i\le x_j,\quad
  x_j+\widehat w_j\le x_i,\quad
  y_i+\widehat h_i\le y_j,\quad
  y_j+\widehat h_j\le y_i.
\]
Ta dùng các biến chứng \(l_{ij}\) và \(u_{ij}\):
\[
  l_{ij}\lor l_{ji}\lor u_{ij}\lor u_{ji},
\]
với quan hệ kéo theo từ từng biến chứng sang bất đẳng thức tương ứng. Các biến
chứng không cần loại trừ lẫn nhau: một vật có thể vừa ở bên trái vừa ở dưới vật
khác.

\begin{designrule}{Quy ước biên}
Tọa độ biểu diễn cạnh của hình chữ nhật, không phải tâm hay ô đã chiếm. Vì thế
dấu \(\le\) cho phép hai vật chạm cạnh. Quy ước này phải được giữ nguyên trong
điều kiện nằm gọn, không chồng lấn và bộ kiểm tra.
\end{designrule}

Hình \ref{fig:packing-layout} dùng một ví dụ tự dựng để minh họa cách tọa độ
thứ tự và quan hệ vị trí được dùng trong SAT cho đóng gói; cách tổ chức quan hệ
này tiếp nối mô hình đóng gói dải của \textcite{soh2010strip}.

\begin{figure}[tb]
\centering
\begin{tikzpicture}[satfig,x=8mm,y=8mm]
  \draw[step=1,Rule!35,very thin] (0,0) grid (8,6);
  \draw[Ink,very thick] (0,0) rectangle (8,6);
  \draw[Blue,fill=PaleBlue,thick] (0,0) rectangle (3,2);
  \node at (1.5,1) {\(A:3\times2\)};
  \draw[Teal,fill=PaleTeal,thick] (3,0) rectangle (5,4);
  \node[rotate=90] at (4,2) {\(B:2\times4\)};
  \draw[Gold,fill=PaleGold,thick] (5,0) rectangle (8,3);
  \node at (6.5,1.5) {\(C:3\times3\)};
  \draw[Blue!65!Teal,fill=SoftGray,thick] (0,4) rectangle (5,6);
  \node at (2.5,5) {\(D:5\times2\)};
  \foreach \t in {0,...,8}
    \node[below,satlabel] at (\t,0) {\(\t\)};
  \foreach \t in {1,...,6}
    \node[left,satlabel] at (0,\t) {\(\t\)};
  \draw[satdim] (0,6.45)--(8,6.45)
    node[midway,above,satlabel]{\(W=8\)};
  \draw[satdim] (8.45,0)--(8.45,6)
    node[midway,right,satlabel]{\(H=6\)};
\end{tikzpicture}
\caption{Một bố trí nguyên có bốn vật. Các cặp chạm cạnh vẫn hợp lệ; chẳng hạn
\(x_A+w_A=x_B=3\), nên biến chứng \(l_{AB}\) có thể nhận giá trị đúng.}
\label{fig:packing-layout}
\end{figure}

\begin{workedexample}{Từ một cặp hình đến literal quan hệ}
Trong \cref{fig:packing-layout}, \(A\) có \((x_A,y_A)=(0,0)\),
\((w_A,h_A)=(3,2)\); \(B\) có \((x_B,y_B)=(3,0)\),
\((w_B,h_B)=(2,4)\). Bốn bất đẳng thức tách là
\[
  0+3\le3,\quad 3+2\le0,\quad 0+2\le0,\quad 0+4\le0.
\]
Chỉ bất đẳng thức đầu đúng, nên \(l_{AB}=1\) là biến chứng tự nhiên. Nếu ép
\(x_B\le2\), miền tọa độ làm \(x_A+w_A\le x_B\) bất khả và mệnh đề liên kết
buộc \(\neg l_{AB}\); phép tuyển bốn quan hệ khi đó phải chọn một hướng khác
hoặc chứng minh cặp hình không thể cùng một bố trí.
\end{workedexample}

\section{Mã hóa thứ tự cho tọa độ}

\subsection{Ngữ nghĩa}

Với \(e\) trong miền ngang, đặt
\[
  p^x_{i,e}=1 \iff x_i\le e.
\]
Tương tự, \(p^y_{i,f}=1\iff y_i\le f\). Chuỗi phải đơn điệu:
\[
  p^x_{i,e}\rightarrow p^x_{i,e+1},
  \qquad
  p^y_{i,f}\rightarrow p^y_{i,f+1}.
\]
Giá trị tọa độ được xác định bởi điểm chuyển đầu tiên từ sai sang đúng. Một
bất đẳng thức như \(x_i+a\le x_j\) trở thành quan hệ kéo theo giữa hai ngưỡng:
\[
  p^x_{i,e}\ \text{và}\ p^x_{j,e+a-1}
\]
được nối theo hướng phù hợp với ngữ nghĩa \(\le\). Lợi ích là một literal
ngưỡng đại diện cho cả nửa miền, thay vì liệt kê từng cặp tọa độ.

\subsection{Hướng đặt và liên kết biểu diễn}

Khi kích thước phụ thuộc \(\rho_i\), quan hệ kéo theo
\[
  l_{ij}\rightarrow x_i+\widehat w_i\le x_j
\]
được tách theo các hướng đặt của \(i\) và \(j\). Các trường hợp không thể nằm
gọn được loại ở bước tiền xử lý. Với vật vuông, đặt
\(\rho_i=0\) để bỏ đối xứng quay.

Khi \(l_{ij}\) đóng vai trò biến chứng cho phép tuyển, mô-đun cần chiều
\[
  l_{ij}\rightarrow x_i+\widehat w_i\le x_j.
\]
Chiều ngược
\[
  x_i+\widehat w_i\le x_j\rightarrow l_{ij}
\]
có thể được bổ sung khi literal quan hệ cần đặc tả hai chiều. Bộ giải mã đọc
trực tiếp tọa độ và hướng đặt, nhờ đó bộ kiểm tra độc lập với lựa chọn liên kết.

\begin{proposition}[Tính đúng của mô-đun không chồng lấn]
Giả sử mỗi biến chứng \(l_{ij},u_{ij}\) kéo theo đúng bất đẳng thức hình học
của nó, và phép tuyển bốn biến chứng được thỏa cho mọi cặp vật trong cùng khung
chứa. Khi đó mọi bố trí giải mã từ mô hình \SAT đều không giao nhau.
\end{proposition}

\begin{proof}
Với một cặp \(i,j\), phép tuyển cho ít nhất một biến chứng đúng. Quan hệ kéo
theo của biến chứng đó cho một bất đẳng thức tách theo trục \(x\) hoặc \(y\). Vì vậy nội
thất hai hình chữ nhật rời nhau.
\end{proof}

\section{Bài toán đóng gói dải hai chiều}

\subsection{Từ tối ưu chiều cao sang kiểm tra tính khả thi}

Dải có chiều rộng cố định \(W\), còn chiều cao \(H\) cần tối thiểu hóa. Với
mỗi \(H\), dựng công thức
\[
  \Phi_{\mathrm{SPP}}(H)
  \iff
  \text{tồn tại bố trí trong }W\times H.
\]
Tính đơn điệu cho phép tìm kiếm trên cận:
\[
  \Phi_{\mathrm{SPP}}(H)=\SAT
  \Longrightarrow
  \Phi_{\mathrm{SPP}}(H')=\SAT
  \quad\forall H'\ge H.
\]
Cận dưới cơ bản là
\[
  LB=\max\left\{
    \max_i \min(h_i,w_i),
    \left\lceil\frac{\sum_iw_ih_i}{W}\right\rceil
  \right\},
\]
trong đó hạng đầu phải điều chỉnh theo hướng đặt thực sự vừa khung. Cận trên
đến từ một phép đặt theo kinh nghiệm.

Chỉ cận diện tích hiếm khi đủ sát. Các phương pháp chính xác còn dùng hàm khả
thi đối ngẫu, phép chiếu một chiều, tính không tương thích và tọa độ tổng tập
con. Trong SAT, mọi cận dưới hợp lệ đều có hai vai trò: thu hẹp số cận phải gọi
và thu hẹp miền tọa độ trước khi sinh mệnh đề. Nén tọa độ đặc biệt quan trọng:
nếu có thể chứng minh một vật chỉ cần bắt đầu tại các tổng tập con khả dĩ,
chuỗi thứ tự nên chạy trên tập mốc đã nén thay
vì mọi số nguyên từ \(0\) đến \(W-w_i\).

\subsection{Ba chế độ giải}

\begin{center}
\captionof{table}{Ba chế độ tối ưu chiều cao bằng SAT/MaxSAT.}
\label{tab:strip-solving-modes}
\begin{tabularx}{\textwidth}{@{}>{\bfseries\sffamily}p{.23\textwidth}YY@{}}
\toprule
Chế độ & Cấu trúc & Đặc điểm\\
\midrule
\SAT độc lập & Dựng lại \(\Phi(H)\) ở mỗi cận &
Đơn giản, mỗi cận có một \CNF riêng\\
\SAT tăng dần & Miền lớn cố định, kích hoạt \(H\) bằng các giả định &
Giữ mệnh đề đã học giữa các cận\\
\MaxSAT có trọng số & Chiều cao được biểu diễn bằng mệnh đề mềm &
Không cần vòng tìm kiếm ngoài\\
\bottomrule
\end{tabularx}
\end{center}

Mã hóa SAT thứ tự cho đóng gói dải đã được nghiên cứu trước đó cùng đối xứng,
quan hệ vị trí và tái sử dụng mệnh đề đã học
\parencite{soh2010strip}. Vì vậy nghiên cứu tình huống của nhóm được định vị
theo hai câu hỏi kế tiếp: phép quay/đối xứng làm đổi mô hình thế nào, và
SAT--MaxSAT--giải tăng dần có giữ lợi thế trên bộ kiểm chuẩn và bộ giải hiện
tại hay không. Phép so sánh hợp lý phải giữ cùng các cận, miền tọa độ và bộ
kiểm tra.

Trong nghiên cứu của \textcite{kieu2025strip}, \SAT không tăng dần với phép
quay và phá đối xứng chứng nhận 26 giá trị tối ưu trên 41 bài toán; phiên bản
tăng dần chứng nhận 21, còn CP chứng nhận 28. Kết quả này cho thấy giải tăng
dần chỉ có lợi khi công thức cơ sở và mệnh đề đã học còn phù hợp sau khi thay
cận. Với đóng gói dải, thay \(H\) có thể làm thay đổi mạnh miền tọa độ và
những xung đột hữu ích.

\begin{keyidea}{Chứng nhận tối ưu}
Một bố trí tại \(H^\star\) là cận trên. Công thức \(\Phi(H^\star-1)\) không
thỏa là chứng nhận cận dưới. Cặp SAT/UNSAT ở hai cận kề nhau mới xác định
\(H^\star\) là giá trị tối ưu.
\end{keyidea}

\section{Bài toán đóng gói vào khung chứa hai chiều}

\subsection{Hai kiến trúc}

Bài toán dùng nhiều khung chứa \(W\times H\) đồng nhất và tối thiểu hóa số
khung. Có hai kiến trúc chính.

\paragraph{Xếp chồng.}
Ghép \(k\) khung thành một dải \(W\times kH\). Tọa độ \(y_i\) đồng thời xác
định khung. Kiến trúc này tái sử dụng đóng gói dải nhưng phải ngăn vật vượt qua
đường biên \(qH\).

\paragraph{Không xếp chồng.}
Dùng biến \(a_{ib}=1\) khi vật \(i\) nằm trong khung \(b\), và tọa độ cục bộ
trong từng khung. Ta có
\[
  \ExactlyOne(a_{i1},\ldots,a_{ik}).
\]
Điều kiện không chồng lấn của cặp \(i,j\) trong khung \(b\) được gác bởi:
\[
  (a_{ib}\land a_{jb})
  \rightarrow
  (l_{ij}\lor l_{ji}\lor u_{ij}\lor u_{ji}).
\]
Biến sử dụng \(q_b\) nối với phép gán bằng \(a_{ib}\rightarrow q_b\).
Đối xứng giữa các khung được giảm bằng
\[
  q_{b+1}\rightarrow q_b.
\]

Các khung đã dùng tạo thành một tiền tố, vì thế hàm mục tiêu \(\sum_bq_b\)
đúng bằng số khung. Có thể thêm phép gán chuẩn cho một số vật đầu, nhưng mọi
ràng buộc đối xứng phải dựa trên phép hoán vị các khung đồng nhất, không dựa
trên một bố trí kinh nghiệm cụ thể.

\begin{figure}[tb]
\centering
\begin{tikzpicture}[satfig,x=7mm,y=7mm]
  \node[satlabel] at (2.5,4.9) {xếp chồng};
  \draw[Ink,thick] (0,0) rectangle (5,4);
  \draw[Rule,dashed] (0,2)--(5,2);
  \node[left,satlabel] at (0,1) {khung \(1\)};
  \node[left,satlabel] at (0,3) {khung \(2\)};
  \draw[Blue,fill=PaleBlue,thick] (.3,.3) rectangle (2.4,1.7);
  \draw[Teal,fill=PaleTeal,thick] (2.6,2.3) rectangle (4.7,3.7);
  \node[satlabel] at (9,4.9) {không xếp chồng};
  \draw[Ink,thick] (6.2,0) rectangle (10.2,2.6);
  \draw[Ink,thick] (10.8,0) rectangle (14.8,2.6);
  \node[below,satlabel] at (8.2,0) {khung \(1\)};
  \node[below,satlabel] at (12.8,0) {khung \(2\)};
  \draw[Blue,fill=PaleBlue,thick] (6.5,.3) rectangle (8.6,1.7);
  \draw[Teal,fill=PaleTeal,thick] (11.1,.3) rectangle (13.2,1.7);
  \node[satstate] at (8.2,3.5) {\(q_1\)};
  \node[satstate] at (12.8,3.5) {\(q_2\)};
  \draw[satdash] (12.25,3.5)--(8.75,3.5)
    node[midway,above,satlabel]{\(q_2\to q_1\)};
\end{tikzpicture}
\caption{Hai kiến trúc cho đóng gói vào khung: xếp chồng dùng một trục \(y\)
toàn cục; không xếp chồng dùng phép gán và tọa độ cục bộ, kèm tiền tố sử dụng.}
\label{fig:stacking-nonstacking}
\end{figure}

Trong mô hình theo vị trí đặt, một cách khác là xây đồ thị xung đột cho mọi cặp
vị trí--khung rồi thêm ràng buộc đúng một theo vật. Cách này tránh biến quan hệ
nhưng có thể tạo số vị trí rất lớn. Mô hình tọa độ không xếp chồng làm ngược
lại: số biến tọa độ nhỏ hơn toàn bộ tập vị trí, đổi lại mỗi cặp vật--khung cần
một phép tuyển có điều kiện. Hai mô hình là hai phía của phép phân tích nhân
tử; đánh giá thực nghiệm nên báo cả số vị trí bị loại ở tiền xử lý và số quan
hệ từng cặp còn lại.

\begin{resultbox}{So sánh kiến trúc}
Nghiên cứu của \textcite{kieu2025bin} so sánh các mô hình \SAT với MIP và CP
trên cùng họ bài toán kiểm chuẩn. Kết quả nhấn mạnh hai yếu tố hơn là một phép
mã hóa đơn lẻ: biến thứ tự giúp biểu diễn bất đẳng thức tọa độ, còn phá đối
xứng cho khung/bản sao quyết định đáng kể kích thước không gian tìm kiếm.
\end{resultbox}

\section{Bài toán cắt tấm một kích thước}

\subsection{Mô hình theo số tờ}

Cho tấm vật liệu kích thước \(W\times H\), nhiều loại vật \(t\), mỗi loại có
nhu cầu \(d_t\). Bài toán tối thiểu hóa số tấm. Nếu tạo từng bản sao riêng,
phép gán và không chồng lấn giống bài toán đóng gói vào khung. Nếu gộp theo
loại, cần thêm ràng buộc đếm để bảo đảm đúng nhu cầu.

Với một cận \(K\), công thức khả thi hỏi liệu có thể đáp ứng đủ nhu cầu trên
\(K\) tấm. Hàm mục tiêu là
\[
  \min\set{K\mid \Phi_{\mathrm{CSP}}(K)\text{ là SAT}}.
\]
Đối xứng xuất hiện ở hai tầng: các tấm đồng nhất và các bản sao cùng loại. Một
cách phá đối xứng bền vững là buộc biến sử dụng tấm theo tiền tố và buộc chỉ
số tấm của các bản sao cùng loại không giảm.

\subsection{SAT và MaxSAT}

Tìm kiếm SAT thay \(K\) ở vòng ngoài. \MaxSAT bộ phận có thể giữ tất cả các
tấm ứng viên, đặt mệnh đề mềm \(\neg q_b\) và tối thiểu
\(\sum_bq_b\). Hai cách dùng chung điều kiện nằm gọn, hướng đặt và không chồng
lấn; chúng chỉ khác ở lớp hàm mục tiêu.

\textcite{kieu2026cutting} cho thấy phép mã hóa gọn với tọa độ thứ tự và đối
xứng theo bản sao/tấm đạt kết quả cạnh tranh trên bộ bài toán cắt tấm một kích
thước. Điểm chuyển giao sang các bài toán đóng gói khác là chuỗi mô-đun
\[
  \texttt{Assignment}
  +\texttt{OrderCoordinate}
  +\texttt{ConditionalNonOverlap}
  +\texttt{UsageCounter}.
\]

\section{Giảm miền và đối xứng}

Trước khi sinh \CNF, bốn phép rút gọn thường có hiệu quả trực tiếp:
\begin{enumerate}
  \item loại hướng đặt không vừa khung;
  \item thu hẹp \(x_i\le W-\widehat w_i\) và
        \(y_i\le H-\widehat h_i\);
  \item loại một hướng tách nếu các cận chứng minh hướng đó bất khả thi;
  \item chuẩn hóa vật/bản sao và khung chứa theo một thứ tự cố định.
\end{enumerate}

Mỗi ràng buộc phá đối xứng được gắn với một nhóm đối xứng cụ thể:
\begin{itemize}
  \item \(\rho_i=0\) cho vật vuông: đối xứng hướng đặt;
  \item \(q_{b+1}\rightarrow q_b\): hoán vị khung/tấm;
  \item thứ tự bản sao cùng loại: hoán vị các bản sao đồng nhất;
  \item neo một vật lớn vào nửa khung chứa: phản xạ/tịnh tiến, nếu phép biến
        đổi giữ miền và hàm mục tiêu.
\end{itemize}
Các nguyên tắc tổng quát về đối xứng được trình bày trong
\textcite{walsh2006symmetry,gent2006symmetry}; phép mã hóa SAT cổ điển cho bài toán đóng gói khung 2D
Packing có rotation có thể xem ở \textcite{soh2010binpacking}.

\section{Tách ảnh hưởng của phép mã hóa, cận và quan hệ trội}

Ma trận thiết kế ghi riêng ảnh hưởng của các thành phần:
\begin{itemize}
  \item cận dưới theo diện tích, phép chiếu và đồ thị con đầy đủ của quan hệ
        không tương thích;
  \item loại hướng đặt và vị trí bị trội;
  \item cố định theo vật lớn hoặc thứ tự vật đồng nhất;
  \item nén tọa độ;
  \item khởi động ấm/cận trên từ phương pháp kinh nghiệm.
\end{itemize}

Ma trận thiết kế tách phần giảm mệnh đề do lọc miền khỏi phần giảm do chuyển
sang CNF. Phép so sánh CP, MIP và SAT dùng chung cách chuẩn hóa bài toán và cận
đầu vào; cận và quan hệ trội được ghi thành mô-đun riêng
\parencite{iori2021packing,clautiaux2007orthogonal}.

\section{Một thư viện mô-đun cho hình học}

\begin{summarybox}
\textbf{OrderCoordinate} cung cấp chuỗi ngưỡng và phép giải mã duy nhất.
\textbf{Orientation} trả về kích thước hiệu dụng.
\textbf{Containment} siết miền theo khung chứa đang chọn.
\textbf{RelativePosition} tạo biến chứng trái/phải/trên/dưới.
\textbf{ConditionalNonOverlap} gác phép tuyển bằng phép gán.
\textbf{UsageCounter} biến số khung chứa được dùng thành hàm mục tiêu.
\end{summarybox}

Thứ tự kết hợp đi từ hướng đặt đến điều kiện nằm gọn, vị trí tương đối, không
chồng lấn được gác bởi phép gán và bộ đếm sử dụng. Bộ kiểm tra đọc vật, hướng
đặt, tọa độ và khung chứa, độc lập với biến phụ.

\section{Bài học thiết kế}

\begin{enumerate}
  \item Chọn một quy ước tọa độ và giữ nó nhất quán trong toàn bộ phép mã hóa.
  \item Dùng biến chứng cho phép tuyển hình học; không buộc biến chứng trở thành
        quan hệ đầy đủ nếu phép giải mã không cần.
  \item Tách mô-đun khả thi khỏi hàm mục tiêu số khung/tấm/chiều cao.
  \item Phá đối xứng theo đúng phép hoán vị hoặc phản xạ của mô hình.
  \item Chứng nhận tối ưu gồm bố trí hợp lệ và một cận dưới tương ứng.
\end{enumerate}

\subsection*{Câu hỏi nghiên cứu}
\begin{enumerate}
  \item Dẫn xuất các mệnh đề thứ tự cho quan hệ kéo theo
        \(l_{ij}\rightarrow x_i+\widehat w_i\le x_j\) ở một hướng đặt cố định.
  \item So sánh số biến của xếp chồng và không xếp chồng trên \(n\) vật,
        \(k\) khung.
  \item Thiết kế ràng buộc phá đối xứng cho ba bản sao cùng loại và chứng minh
        mỗi quỹ đạo giữ một
        đại diện.
\end{enumerate}
